\section{Modelo de dominio}

\subsection{Descripcion}

Este documento presenta el dominio funcional del sistema y sus agrupaciones conceptuales principales. Su objetivo es expresar el lenguaje del negocio antes de cualquier detalle de tablas, endpoints o decisiones fisicas de persistencia.

\subsection{Alcance}

Aplica al dominio central de la plataforma de quiniela.

\subsection{Definiciones}

\begin{itemize}
  \item \textbf{Dominio}: conjunto de reglas y conceptos que definen como funciona la quiniela.
  \item \textbf{Subdominio}: area funcional con responsabilidad delimitada dentro del sistema.
\end{itemize}

\subsection{Reglas}

Subdominios principales:

\begin{itemize}
  \item `Identidad y acceso`: registro, autenticacion y estado de usuario.
  \item `Activacion y pagos`: comprobante, revision administrativa y habilitacion posterior del usuario.
  \item `Torneo y calendario`: torneo, equipos, partidos, fases y kickoff.
  \item `Picks`: predicciones por partido y predicciones especiales de torneo.
  \item `Resultados`: recepcion de resultados oficiales y correcciones.
  \item `Scoring`: calculo de puntos, multiplicadores y desempates.
  \item `Leaderboard`: posicionamiento general y filtrado.
  \item `Administracion`: acciones operativas de la JD.
  \item `Auditoria`: trazabilidad de acciones y cambios.
  \item `Notificaciones`: eventos comunicables por correo.
\end{itemize}

Reglas de separacion conceptual:

\begin{itemize}
  \item El estado funcional del usuario pertenece a `Identidad y acceso`.
  \item La revision del comprobante pertenece a `Activacion y pagos` y no crea por si misma un nuevo estado de usuario distinto de \texttt{PENDIENTE}.
  \item Las predicciones especiales pertenecen al subdominio `Picks`, aunque su ventana de captura y su evaluacion sean distintas de los picks por partido.
  \item `Resultados` mantiene una unica verdad vigente por partido para efectos de scoring, aunque pueda existir historia de correcciones y overrides.
  \item `Leaderboard` es un subdominio derivado: no define reglas propias de juego, solo expresa el resultado acumulado del scoring y sus desempates.
\end{itemize}

Relaciones de dominio:

\begin{itemize}
  \item Un usuario participa en un torneo mediante picks por partido y predicciones especiales.
  \item Un comprobante de pago puede sustentar la activacion manual de un usuario, pero no equivale por si mismo a habilitacion competitiva.
  \item Un partido pertenece a una fase del torneo y usa un kickoff oficial vigente para su lock.
  \item Un resultado oficial vigente alimenta el scoring de todos los picks asociados al partido.
  \item El scoring consolida puntaje por partido y puntaje de predicciones especiales para actualizar el leaderboard.
  \item Acciones administrativas, correcciones y recalculos generan auditoria.
\end{itemize}

\subsection{Casos borde}

\begin{itemize}
  \item El dominio distingue entre resultado oficial y override manual, pero ambos deben converger a una sola verdad vigente.
  \item El dominio debe soportar ausencia de pick sin convertirla en un pick implicito.
  \item El dominio debe soportar partidos postergados, suspendidos, cancelados o corregidos sin romper la trazabilidad de scoring.
\end{itemize}

\subsection{Invariantes}

\begin{itemize}
  \item Las reglas del juego pertenecen al dominio y no a UI ni a integraciones.
  \item Un resultado oficial vigente debe ser unico por partido a efectos de scoring.
  \item Ninguna correccion de resultado puede eliminar la historia operativa previa.
  \item La ausencia de pick es un resultado funcional valido del dominio, no un valor por defecto fabricado por el sistema.
\end{itemize}

\subsection{Preguntas abiertas}

\begin{itemize}
  \item Confirmar si el modelo conceptual necesita distinguir explicitamente entre puntaje calculado y snapshot persistido de leaderboard cuando se documente el modelo derivado.
\end{itemize}
